\documentclass[12pt]{article}
\usepackage[utf8]{inputenc}
\usepackage[T1]{fontenc}
\usepackage{hyperref}
\title{Tenis Language Resources for K-12 Curriculum Development}
\author{Compiled by Claude Sonnet 5 (claude-sonnet-5) for Andr\'as Kornai}
\date{}
\begin{document}
\maketitle

\section{Introduction}
Tenis (Tench; ISO 639-3 tns) is an Oceanic language of the St Matthias group. It is spoken on Tench Island, New Ireland Province, Papua New Guinea. Counts differ: about 30 speakers (Wikipedia, 2000) and 49 reported in 1972. UNESCO lists it as critically endangered. Speakers are highly bilingual in Mussau-Emira, which puts pressure on Tenis.

\section{Primary Educational Resources}
None found.

\section{Linguistic and Reference Materials}
\begin{itemize}
\item Wikipedia: \url{https://en.wikipedia.org/wiki/Tenis_language}
\item Endangered Languages Project entry: \url{https://www.endangeredlanguages.com/lang/2693}
\item The Avondale repository record ``Mussau, Emirau, and Tench Islands, Papua New Guinea'' appeared in results. I did not open it, so its relevance to Tenis is unverified: \url{https://research.avondale.edu.au/handle/1897/19491117}
\end{itemize}

\section{Religious and Community Resources}
None found.

\section{Digital and Audio Resources}
None found.

\section{Curriculum Development Context}
No grammar, dictionary or teaching materials were found. Documentation would have to come before any curriculum, and the elderly speakers are few. This dossier is short because little was found.
\end{document}